\section*{Введение}
\addcontentsline{toc}{section}{Введение}

В лабораторной работе решалась задача прогнозирования интегрального показателя
субъективного благополучия респондента --- \textit{ощущаемого счастья}. Данные
содержат две группы признаков: причины, доступные для всех строк, и состояния,
которые заполнены только для наблюдений с известным классом счастья. Для 600
наблюдений целевой признак равен \code{Неизвестно}, а признаки состояния в этих
строках полностью отсутствуют.

Из-за этого прямое обучение классификатора на реальных состояниях некорректно:
на known-части оно дает завышенную оценку качества, но не может быть применено к
unknown-строкам. Поэтому основной подход в работе был двухэтапным:
\begin{center}
\code{причины $\rightarrow$ восстановленные состояния $\rightarrow$ счастье}.
\end{center}

\textbf{Цель работы} --- построить воспроизводимую модель для прогноза счастья,
достигающую требуемого уровня micro precision не ниже $0{,}6$, и объяснить,
какие признаки и ограничения данных определяют качество прогноза.

Для достижения цели были поставлены следующие \textbf{задачи}:
\begin{enumerate}[noitemsep]
    \item удалить мультиколлинеарные причины с помощью VIF-фильтра;
    \item отобрать состояния, наиболее связанные с рангом счастья;
    \item восстановить выбранные состояния по причинам и сфере занятости;
    \item проверить нелинейные преобразования причин и их влияние на $R^2$;
    \item изучить роль сообществ как источника групповой структуры данных;
    \item обучить логистическую регрессию и получить прогноз для unknown-строк.
\end{enumerate}

\newpage
